\scanpage{45}
Notons que pour deux espaces vectoriels $E,F$, on a un isomorphisme canonique
\[
(2)\qquad\overline{E\otimes F}=\overline E\otimes\overline F.
\]
En effet, l'application $(\overline x,\overline y)\longrightarrow\unsure{\overline{x\otimes y}|\overline{\overline x\otimes\overline y}}$ de $\overline E\unsure{\times|\otimes}\overline F$ dans $\overline{E\otimes F}$ est manifestement bilinéaire, donc se prolonge en une application linéaire $\overline E\otimes\overline F\longrightarrow\overline{E\otimes F}$.
On vérifie alors aussitôt, au moyen des bases de $E$ et $F$ par exemple, que \unsure{c'est là} un isomorphisme du premier espace \uline{sur} le second.

\deleted{Supposons} Rappelons qu'on appelle espace \uline{préhilbertien} un espace vectoriel $E$, muni d'une forme sesquilinéaire sur $E\times E$, \deleted{\unsure{à valeurs dans C}} \inserted{dont la} valeur pour les éléments $x,y$ est notée $\langle x,y\rangle$, satisfaisant à la condition $\langle x,x\rangle\geq0$ pour tout $x$. Alors,
\deleted{\unsure{Nous obtenons ainsi une} \unsure{fonction} \illegible{3} $E\unsure{\otimes}$}
$\|x\|=\sqrt{\langle x,x\rangle}$ est une semi-norme sur $E$, et $E$ est dit un \uline{espace de Hilbert} si $E$ est séparé et complet pour cette norme.
\deleted{\illegible{5}}\deleted{\unsure{Séparé hilbertien}}
Notons \unsure{uniquement}.....

\margin{\deleted{Si à tout $y\in E$ on fait correspondre la forme linéaire $x\longrightarrow(x,y)$ sur $E$, on définit une application \deleted{\unsure{linéaire}} \inserted{\unsure{antilinéaire}} de $E$ sur le dual $E'$, conservant la norme. Par suite, on a un isomorphisme \uline{\unsure{canonique}} \unsure{d'espaces} \illegible{1}
\[
(3)\qquad E'=\overline E\qquad(E\text{ Hilbert}).
\]}}

\unsure{Soient maintenant} $E$ et $F$ deux espaces \deleted{préhilbertiens} \inserted{de Hilbert}, on \unsure{mettra} sur $E\otimes F$ une structure \deleted{préhilbertienne} \inserted{d'espace préhilbertien séparé} \inserted{canonique}.
Pour ceci, on doit définir une forme linéaire sur $(E\otimes F)\otimes\overline{E\otimes F}$ ; \deleted{\unsure{donc}} en vertu de (2), \deleted{\unsure{une forme}} cet espace est isomorphe canoniquement à $(E\otimes\overline E)\otimes(F\otimes\overline F)$ ; il faut donc définir une forme linéaire sur ce dernier, \unsure{on en connaît une}
\margin{Introduire sur $E\otimes F\subset L(\overline E,F)$ la forme \uline{$\operatorname{Tr}(v^*u)$} directement}
\mspage{46}
